% GENERATED FILE. Edit canonical lessons, then run npm run generate:books.
% canonical-source-hash: fnv1a64:205ce8e16a26a89e
% canonical-lessons: FA-C531-mostanad, FA-C531-kartun, FA-C531-zirnevis, FA-C531-duble, FA-C531-kanal

\chapter{On Screen}
\label{ch:fa-on-screen}
\hlchaptermodality{531}

\begin{chapteropening}
\textbf{By the end of this chapter:} \emph{I can talk about what I watch on screen.}

Complete the last lesson of chapter 531: I can talk about what I watch on screen.
\end{chapteropening}

\section[\texorpdfstring{\fa{مستند} (mostanad)}{مستند (mostanad)}]{\fa{مستند} (mostanad) — a documentary}

\label{lesson:FA-C531-mostanad}



\begin{quote}
Before the new one: say the Persian for tragic, then the Persian for funny.
\end{quote}

\subsection*{You'll want to know: \fa{مستند}}
\textbf{\fa{مستند}} (\emph{mostanad}) — \textquotedblleft{}a documentary\textquotedblright{}.

\begin{quote}
\fa{او} \fa{مستندهای} \fa{طبیعت} \fa{را} \fa{دوست} \fa{دارد}.
\end{quote}

\begin{quote}
\emph{u mostanadhâ-ye tabi'at râ dust dârad.}
\end{quote}

\begin{quote}
He likes nature documentaries.
\end{quote}

Also an adjective 'documented': \emph{film-e mostanad}.

\begin{cousinweb}
\emph{mostanad} is a loan, usually traced to Arabic \emph{mustanad}, 'a document, something relied on'. EtymDB has no entry for it.
\end{cousinweb}

\begin{grammarlens}[title={What you've built}]
One more word for telling, explaining and describing at B1.
\end{grammarlens}

\subsection*{Your turn}
\begin{itemize}
\raggedright
  \item \emph{Say it:} \emph{mostanad}
  \item \emph{Say it:} \emph{mostanad}, once more
  \item \emph{Say it:} a sentence of your own with \emph{mostanad}
  \item \emph{Recall:} say the Persian for tragic, then the Persian for funny, then say \emph{mostanad} again
\end{itemize}

\begin{tcolorbox}[breakable,colback=teal!4,colframe=teal!35!black,title={Before you move on}]
What does \fa{مستند} mean? (\textquotedblleft{}A documentary\textquotedblright{}.) Say it once more.
\end{tcolorbox}
\section[\texorpdfstring{\fa{کارتون} (kârtun)}{کارتون (kârtun)}]{\fa{کارتون} (kârtun) — a cartoon (animated film)}

\label{lesson:FA-C531-kartun}



\begin{quote}
Before the new one: say the Persian for funny, then the Persian for a documentary.
\end{quote}

\subsection*{You'll want to know: \fa{کارتون}}
\textbf{\fa{کارتون}} (\emph{kârtun}) — \textquotedblleft{}a cartoon (animated film)\textquotedblright{}.

\begin{quote}
\fa{بچه} \fa{ها} \fa{صبحها} \fa{کارتون} \fa{تماشا} \fa{میکنند}.
\end{quote}

\begin{quote}
\emph{bachchehâ sobhhâ kârtun tamâshâ mikonand.}
\end{quote}

\begin{quote}
The children watch cartoons in the mornings.
\end{quote}

Not \emph{kârton} 'cardboard box', which is spelled \fa{کارتن}.

\begin{cousinweb}
\emph{kârtun} was borrowed from English \emph{cartoon} (EtymDB).
\end{cousinweb}

\begin{grammarlens}[title={What you've built}]
One more word for telling, explaining and describing at B1.
\end{grammarlens}

\subsection*{Your turn}
\begin{itemize}
\raggedright
  \item \emph{Say it:} \emph{kârtun}
  \item \emph{Say it:} \emph{kârtun}, once more
  \item \emph{Say it:} a sentence of your own with \emph{kârtun}
  \item \emph{Recall:} say the Persian for funny, then the Persian for a documentary, then say \emph{kârtun} again
\end{itemize}

\begin{tcolorbox}[breakable,colback=teal!4,colframe=teal!35!black,title={Before you move on}]
What does \fa{کارتون} mean? (\textquotedblleft{}A cartoon (animated film)\textquotedblright{}.) Say it once more.
\end{tcolorbox}
\section[\texorpdfstring{\fa{زیرنویس} (zirnevis)}{زیرنویس (zirnevis)}]{\fa{زیرنویس} (zirnevis) — subtitles}

\label{lesson:FA-C531-zirnevis}



\begin{quote}
Before the new one: say the Persian for a documentary, then the Persian for a cartoon.
\end{quote}

\subsection*{You'll want to know: \fa{زیرنویس}}
\textbf{\fa{زیرنویس}} (\emph{zirnevis}) — \textquotedblleft{}subtitles\textquotedblright{}.

\begin{quote}
\fa{من} \fa{فیلمهای} \fa{خارجی} \fa{را} \fa{با} \fa{زیرنویس} \fa{میبینم}.
\end{quote}

\begin{quote}
\emph{man filmhâ-ye khâreji râ bâ zirnevis mibinam.}
\end{quote}

\begin{quote}
I watch foreign films with subtitles.
\end{quote}

\emph{zir} 'under' + \emph{nevis} 'writing'.

\begin{cousinweb}
\emph{zirnevis} is built on native Persian \emph{zir}, 'under', usually traced to Middle Persian \emph{azēr}, and \emph{nevis}, the present stem of \emph{neveshtan}, 'to write': under-writing.
\end{cousinweb}

\begin{grammarlens}[title={What you've built}]
One more word for telling, explaining and describing at B1.
\end{grammarlens}

\subsection*{Your turn}
\begin{itemize}
\raggedright
  \item \emph{Say it:} \emph{zirnevis}
  \item \emph{Say it:} \emph{zirnevis}, once more
  \item \emph{Say it:} a sentence of your own with \emph{zirnevis}
  \item \emph{Recall:} say the Persian for a documentary, then the Persian for a cartoon, then say \emph{zirnevis} again
\end{itemize}

\begin{tcolorbox}[breakable,colback=teal!4,colframe=teal!35!black,title={Before you move on}]
What does \fa{زیرنویس} mean? (\textquotedblleft{}Subtitles\textquotedblright{}.) Say it once more.
\end{tcolorbox}
\section[\texorpdfstring{\fa{دوبله} (duble)}{دوبله (duble)}]{\fa{دوبله} (duble) — dubbed; dubbing}

\label{lesson:FA-C531-duble}



\begin{quote}
Before the new one: say the Persian for a cartoon, then the Persian for subtitles.
\end{quote}

\subsection*{You'll want to know: \fa{دوبله}}
\textbf{\fa{دوبله}} (\emph{duble}) — \textquotedblleft{}dubbed; dubbing\textquotedblright{}.

\begin{quote}
\fa{پدربزرگ} \fa{فیلمهای} \fa{دوبله} \fa{را} \fa{بیشتر} \fa{دوست} \fa{دارد}.
\end{quote}

\begin{quote}
\emph{pedarbozorg filmhâ-ye duble râ bishtar dust dârad.}
\end{quote}

\begin{quote}
Grandfather prefers dubbed films.
\end{quote}

\emph{duble kardan} — to dub.

\begin{cousinweb}
\emph{duble} is a modern loan, usually traced to French \emph{doublé}, 'doubled'. EtymDB has no entry for it.
\end{cousinweb}

\begin{grammarlens}[title={What you've built}]
One more word for telling, explaining and describing at B1.
\end{grammarlens}

\subsection*{Your turn}
\begin{itemize}
\raggedright
  \item \emph{Say it:} \emph{duble}
  \item \emph{Say it:} \emph{duble}, once more
  \item \emph{Say it:} a sentence of your own with \emph{duble}
  \item \emph{Recall:} say the Persian for a cartoon, then the Persian for subtitles, then say \emph{duble} again
\end{itemize}

\begin{tcolorbox}[breakable,colback=teal!4,colframe=teal!35!black,title={Before you move on}]
What does \fa{دوبله} mean? (\textquotedblleft{}Dubbed; dubbing\textquotedblright{}.) Say it once more.
\end{tcolorbox}
\section[\texorpdfstring{\fa{کانال} (kânâl)}{کانال (kânâl)}]{\fa{کانال} (kânâl) — a TV channel; a canal}

\label{lesson:FA-C531-kanal}



\begin{quote}
Before the new one: say the Persian for subtitles, then the Persian for dubbed.
\end{quote}

\subsection*{You'll want to know: \fa{کانال}}
\textbf{\fa{کانال}} (\emph{kânâl}) — \textquotedblleft{}a TV channel; a canal\textquotedblright{}.

\begin{quote}
\fa{او} \fa{کانال} \fa{تلویزیون} \fa{را} \fa{عوض} \fa{کرد}.
\end{quote}

\begin{quote}
\emph{u kânâl-e televiziyon râ 'avaz kard.}
\end{quote}

\begin{quote}
He changed the TV channel.
\end{quote}

Also a waterway, and a channel online.

\begin{cousinweb}
\emph{kânâl} was borrowed from French \emph{canal} (EtymDB).
\end{cousinweb}

\begin{grammarlens}[title={What you've built}]
One more word for telling, explaining and describing at B1.
\end{grammarlens}

\subsection*{Your turn}
\begin{itemize}
\raggedright
  \item \emph{Say it:} \emph{kânâl}
  \item \emph{Say it:} \emph{kânâl}, once more
  \item \emph{Say it:} a sentence of your own with \emph{kânâl}
  \item \emph{Recall:} say the Persian for subtitles, then the Persian for dubbed, then say \emph{kânâl} again
\end{itemize}

\begin{tcolorbox}[breakable,colback=teal!4,colframe=teal!35!black,title={Before you move on}]
What does \fa{کانال} mean? (\textquotedblleft{}A TV channel; a canal\textquotedblright{}.) Say it once more.
\end{tcolorbox}
